% 論文節の日本語訳（読み物 PDF）用 preamble
% lualatex + jlreq
\documentclass[12pt,paper=a4]{jlreq}
\usepackage{luatexja-fontspec}
\usepackage{amsmath,amssymb,bm}
\usepackage{graphicx}
\usepackage{hyperref}
\usepackage{enumitem}
\usepackage{booktabs}
\usepackage{microtype}
\usepackage{xcolor}
\usepackage{caption}
\usepackage{tcolorbox}
\tcbuselibrary{breakable,skins}

\hypersetup{
  colorlinks=true,
  linkcolor=blue!60!black,
  urlcolor=blue!60!black,
  citecolor=blue!60!black,
}

\graphicspath{{images/}{./}}

\setlength{\parindent}{1em}
\setlength{\parskip}{0.55em}
\setlist[itemize]{leftmargin=1.6em,itemsep=0.25em,topsep=0.35em}
\setlist[enumerate]{leftmargin=1.8em,itemsep=0.25em,topsep=0.35em}

\captionsetup{
  font=small,
  labelfont=bf,
  skip=0.4em,
  justification=raggedright,
  singlelinecheck=false,
}

\newcommand{\vx}{\bm{x}}
\newcommand{\vu}{\bm{u}}
\newcommand{\veps}{\bm{\epsilon}}
\newcommand{\vw}{\bm{w}}

\DeclareMathOperator*{\argmax}{arg\,max}
\DeclareMathOperator*{\argmin}{arg\,min}
\DeclareMathOperator{\Tr}{Tr}
\DeclareMathOperator{\diag}{diag}

% 命題・定理ボックス（ページ跨ぎ時も切れ目を抑える）
\newtcolorbox{propositionbox}[1][]{%
  enhanced,
  breakable,
  colback=blue!4!white,
  colframe=blue!45!black,
  coltitle=white,
  fonttitle=\bfseries\small,
  title={#1},
  attach boxed title to top left={yshift=-2mm, xshift=3mm},
  boxed title style={
    colback=blue!45!black,
    colframe=blue!45!black,
    boxrule=0pt,
    arc=1pt,
  },
  boxrule=0.5pt,
  arc=1.5pt,
  left=5pt, right=5pt, top=8pt, bottom=5pt,
  before skip=0.9em,
  after skip=0.9em,
  pad at break=3mm,
}

% 図はフロートせず本文位置に置く（分断を防ぐ）
\newcommand{\translationfigure}[2]{%
  \begin{center}
    \includegraphics[width=0.88\linewidth]{#1}\\[0.35em]
    {\small #2\par}
  \end{center}
  \vspace{0.4em}%
}

\begin{document}
{\color{gray}\footnotesize 出典: §2.3\par}
\vspace{0.2em}

\section*{2.3 結びの言葉}


本章では、変分的レンズを通じて拡散モデルの起源を辿った。まず変分オートエンコーダ（VAE）——エビデンス下界（ELBO）を通じてデータと構造化された潜在空間の間の確率的写像を学習する基礎的生成モデル——から始めた。階層的 VAE（HVAE）が潜在層を積み重ねることでこの考えを拡張し、漸進的な粗いものから細かいものへの生成という強力な概念を導入したことを見た。しかしこれらのモデルは、学習の安定性とサンプル品質の課題に直面する。

次に、デノイジング拡散確率モデル（DDPMs）を、この変分枠組み内の決定的な進化として枠づけた。エンコーダを漸進的なノイズ付加過程に固定し、逆デノイジングステップのみを学習することで、DDPM は HVAE の学習不安定性を優雅に回避する。決定的に、DDPM も対数尤度の変分下界を最大化することで学習され、その学習目的は一連の単純なデノイジングタスクに分解することを示した。この扱いやすさは、扱い不能な周辺目的を扱い可能な条件付き目的に変換する強力な条件付け戦略によって可能になる——拡散モデルにおける繰り返し現れるテーマである。

この変分枠組みは DDPM に対する完全で強力な基盤を提供するが、それらを理解する唯一の方法ではない。代替的で同様に根本的な視点が、エネルギーベースモデリングの原理から現れる。次章では、このスコアベースの視点を探る：

\begin{enumerate}
\item 焦点をデノイジング遷移確率 $p_\phi(x_{i-1}\mid x_i)$ の学習から、データの対数密度の勾配、すなわちスコア関数の直接学習へ移す。
\end{enumerate}


\begin{enumerate}
\item EBM に起源をもつこのアプローチが、どのようにノイズ条件付きスコアネットワーク（Noise Conditional Score Networks; NCSN）を生じさせ、DDPM で学習されるノイズ予測（$\epsilon$-prediction）とスコア関数そのものの間の深い数学的同値を明らかにするかを見る。
\end{enumerate}


この代替的視点は新しい洞察を与えるだけでなく、後に発展させる拡散モデルの統一的連続時間枠組みのためのもう一つの要石としても機能する。

---

\section*{第3章への導入}


\subsection*{スコアベース視点：EBM から NCSN へ}


これまでの章では、拡散モデルをその変分的ルーツまで辿り、それらが VAE の枠組み内でどのように生じるかを示した。ここで第二の、同様に根本的な視点に転じる：エネルギーベースモデル（Energy-Based Models; EBMs）（Ackley et al., 1985; LeCun et al., 2006）である。EBM は、データ上で低く他では高いエネルギー景観によって分布を表現する。サンプリングは典型的に Langevin 力学に依存し、この景観の勾配に従ってサンプルを高密度領域へ動かす。この勾配場はスコアとして知られ、より高い確率の方向を指す。

中心的な観察は、生成にはスコアを知れば十分であることである：それは扱い不能な正規化定数を計算することなく、サンプルを尤もらしい領域へ動かす。スコアベース拡散モデルはこの考えの上に直接構築される。クリーンなデータ分布だけに焦点を当てる代わりに、スコアが近似しやすい一連のガウスノイズ摂動分布を考える。これらのスコアを学習すると、ノイズの多いサンプルをステップごとにデータへ導くベクトル場の族が得られ、生成を漸進的デノイジングに変える。

本章の後には次ができるようになる：

\begin{itemize}
\item スコアベースモデリングの中核アイデア——スコア関数とは何か、およびそれが Langevin 力学とともに生成にどう使えるか——を理解する；
\item デノイジングスコアマッチングが条件付けトリックを通じて単一ノイズレベルでスコアをどう学習するか、および Tweedie の公式がスコア予測をデノイジングにどう接続するかを理解する；
\item ノイズ条件付きスコアネットワーク（NCSN）がデノイジングスコアマッチングを複数ノイズレベルへどう拡張するか、なぜ複数ノイズレベルが有用か、および NCSN のスコア予測が DDPM のノイズ予測とどう関係するかを理解する。
\end{itemize}


\textbf{読者ルート。} エネルギーベースモデルに既に親しんでいる読者は第3.1節を流し読みし、第3.2.1節から始めてよい。
\end{document}
